\section{问题建模}

\subsection{任务场景与海洋再分析数据}
本文研究自主水下航行器在强流海区携带陈旧粗分辨率海流图执行绕障航行的二元路径决策问题，参数汇总于表\ref{tab:settings}。研究海区取台湾以东至东海的黑潮区（23°--31°N、121°--129°E；第\ref{sec:scope}节另在湾流区与阿古拉斯回流区重复，并改带卫星测高地转流图重复），采用HYCOM全球再分析（GLBy0.08，实验93.0\cite{chassignet2007hycom}）10米水深处的1/12°水平流速场，2019年1月1日至2024年8月29日每3天一张快照共689张；其中2019--2021年作为开发年份（710个任务），2022--2024年作为独立测试年份（623个任务）。每张快照生成2个任务，起终点相距100--140千米、离岸不少于20千米，中点设半径15千米的圆形禁航区，只保留粗图上任务区平均流速不低于0.3米/秒的强流任务。航行器下水时携带3天前生成的1°粗图（格内平均），当天1/12°场作为真实海流。

\begin{table}[htbp]
\centering
\caption{粗分辨率海流图下自主水下航行器航路决策问题的实验设置与参数。}
\label{tab:settings}
\small
\begin{tabular}{@{}p{0.16\linewidth}p{0.80\linewidth}@{}}
\toprule
项目 & 设置 \\
\midrule
海流数据 & HYCOM 全球海洋再分析（GLBy0.08，实验 93.0），水平分辨率 1/12°，取 10 米深度；23°--31°N、121°--129°E（台湾以东至东海，含黑潮）；2019-01-01 至 2024-08-29 每 3 天 12:00 时一张，共 689 张 \\
任务 & 每张快照抽 2 个：起终点相距 100--140 千米、离岸不少于 20 千米，两点中间设半径 15 千米的禁航区，可从左侧或右侧绕过；只取粗图上任务区平均流速不低于 0.3 米/秒的强流任务；开发 710 个（2019--2021 年），测试 623 个（2022--2024 年） \\
地图与真实海流 & 航行器带 3 天前的 1° 粗分辨率海流图（格内平均）；真实海流为当天 1/12° 场 \\
航行器 & 对水速度 1.0 米/秒；选定一侧后，按真实海流在该侧以最快方式航行 \\
航线求解 & 网格最短路：粗算 8 千米网格、8 个方向；标准 4 千米、16 个方向；精算 2 千米、32 个方向（计算量之比 1 : 8 : 64）；顶流过强的边不可通行 \\
合理海流 & 粗图加上由流函数生成的无散误差模态的组合，各点误差大小与相关长度（76 千米）由开发年份标定，组合系数服从标准正态；半径 $r=1.5$（开发年份中包住真实两侧时间差的任务占 93\%） \\
反例检验 & 在半径 $r$ 的球内对两侧时间差做投影梯度上升（梯度沿最快航线累加）；两个方向、半径 1 与 2 各搜一次，其他半径线性插值，每个任务约 71 次标准网格求解 \\
观测方式 & 出发前原地停留 2 小时测 12 次；沿所选一侧航行 30 或 60 千米并连续测流后可换边；绕到反例与地图差别最大处（前方 40 千米内）去测；上浮取当天 1° 或 0.5° 海流图（每次 1 小时）；测流噪声 0.05 米/秒 \\
门控 & 最坏后悔超过 8 小时（开发年份选定）时取保守一侧或上浮取图，否则照地图走 \\
评价 & 损失 = 实际航时 $-$ 更快一侧的航时（小时）；配对差值的 95\% 置信区间按天整组自助重抽样 2000 次 \\
\bottomrule
\end{tabular}

\end{table}

\subsection{航线求解与计算预算}
给定水平海流场 $\boldsymbol{u}(\boldsymbol{x})$ 与航行器恒定对水速度 $V = 1.0$ 米/秒，选定绕行一侧 $s \in \{\text{L}, \text{R}\}$（左侧或右侧）后的最快航线由有向图最短路算法求取\cite{dijkstra1959note}。对长度为 $\ell_e$、单位方向为 $\boldsymbol{d}_e$ 的有向边 $e$，以边两端节点海流的平均作为该边的海流，记沿边分量为 $u_{\parallel} = \boldsymbol{u} \cdot \boldsymbol{d}_e$、垂直边分量为 $u_{\perp} = \|\boldsymbol{u} - u_{\parallel}\boldsymbol{d}_e\|$。航行器抵消横流后的对地速度 $v_{\text{g}}(e)$ 与通行时间 $t(e)$ 为：
\begin{equation}
v_{\text{g}}(e) = u_{\parallel} + \sqrt{V^2 - u_{\perp}^2}, \qquad t(e) = \frac{\ell_e}{v_{\text{g}}(e)}.
\label{eq:edge_time}
\end{equation}
当 $|u_{\perp}| \ge V$ 或 $v_{\text{g}}(e) \le 0$ 时该边不可通行。在严格位于禁航区左侧或右侧的子图上求最短路，得两侧最优航时 $T_{\text{L}}(\boldsymbol{u})$ 与 $T_{\text{R}}(\boldsymbol{u})$ 及两侧时间差：
\begin{equation}
\Delta T(\boldsymbol{u}) = T_{\text{L}}(\boldsymbol{u}) - T_{\text{R}}(\boldsymbol{u}),
\label{eq:delta_t}
\end{equation}
其中 $\Delta T(\boldsymbol{u}) < 0$ 表示左侧快。本文设三档求解计算量：粗算（8千米8向）、标准（4千米16向）与精算（2千米32向），按边数计相对计算量为 $1:8:64$。选定一侧后，假设航行器按真实海流 $\boldsymbol{u}_{\text{true}}$ 在该侧以最优方式航行（测试年份更快一侧平均航时39.2小时），从而聚焦于选左侧还是右侧这一二元决策。

\subsection{合理海流集合的参数化建模}
由于3天前的1°粗图 $\boldsymbol{u}_0(\boldsymbol{x})$ 缺失了中尺度切变与流轴演变，我们用流函数生成无散的误差模态，并按多年数据标定的误差大小逐点缩放，构造叠加在粗图之上的合理海流 $\boldsymbol{u}(\boldsymbol{x}; \boldsymbol{z})$：
\begin{equation}
\boldsymbol{u}(\boldsymbol{x}; \boldsymbol{z}) = \boldsymbol{u}_0(\boldsymbol{x}) + \sum_{k=1}^{K} z_k \boldsymbol{m}_k(\boldsymbol{x}), \qquad \boldsymbol{m}_k(\boldsymbol{x}) = \mathbf{S}(\boldsymbol{x})\, \nabla^{\perp}\phi_k(\boldsymbol{x}),
\label{eq:plausible_field}
\end{equation}
其中 $\nabla^{\perp} = (-\partial_y, \partial_x)$，$\{\phi_k\}_{k=1}^{K}$ 为相关长度 $L_{\text{corr}}$ 的高斯相关流函数的前 $K$ 个主模态（保留95\%方差），因而 $\nabla^{\perp}\phi_k$ 处处无散；$\mathbf{S}(\boldsymbol{x})$ 为对角缩放矩阵，使两个流速分量在各点的方差等于开发数据标定的误差方差（缩放后近似无散）；$\boldsymbol{z} = (z_1, \dots, z_K)^{\top}$ 为先验服从标准正态分布的组合系数。给定半径 $r > 0$，合理集合定义为
\begin{equation}
\mathcal{U}(r) = \left\{ \boldsymbol{u}(\boldsymbol{x}; \boldsymbol{z}) \;\middle|\; \|\boldsymbol{z}\|_2 \le r \right\}.
\label{eq:plausible_set}
\end{equation}
半径 $r$ 由开发年份标定（第5.3节），使 $\mathcal{U}(r)$ 恰好包住大多数真实海洋的两侧时间差。

\subsection{可选动作、观测代价与评价指标}
航行器面对粗图 $\boldsymbol{u}_0$ 可在五类动作中选择：（1）同一张粗图上加算：用粗算、标准或照地图走（精算）选更快一侧出发；（2）改变决策准则：在32个合理海流上按期望时间选边，或采取保守（最坏后悔最小）准则；（3）就地取新信息：原地停留测（2小时测12次）、沿途测30千米、沿途测60千米后再决定是否换边，或绕行定点测（去前方40千米内反例与地图差别最大处测量；测流噪声0.05米/秒；测得的海流按线性高斯更新修正模态系数，观测误差在开发年份比较0.08、0.15、0.3米/秒后取0.15米/秒）；（4）上浮取新图：每次耗时1小时接收当天1°图或0.5°图；（5）本文反例检验门控：仅在最坏后悔超阈值时触发门控后取保守或门控后上浮取图。定义任务航时“损失”为实际总耗时 $T_{\text{total}}$ 减去真实更快一侧的航时：
\begin{equation}
\mathcal{L} = T_{\text{total}} - \min\left\{ T_{\text{L}}(\boldsymbol{u}_{\text{true}}),\, T_{\text{R}}(\boldsymbol{u}_{\text{true}}) \right\}.
\label{eq:loss}
\end{equation}

\subsection{问题剖析：为何常规计算与盲目观测难以奏效}
选对绕行一侧依赖禁航区两侧数十千米尺度的海流切变，而3天前1°粗图抹平了这一信息；机载计算只能锐化粗图自身的几何推论，却无法创造缺失的物理信息。这带来两个结构性性质：

第一，粗图上的时间差模糊性源于“数值歧义”与“环境歧义”两类不同机制。当 $|\Delta T(\boldsymbol{u}_0)|$ 很小时，若仅因8千米粗算网格折线失真，换成2千米精算即可看清粗图倾向；但若精算选出某侧 $s_0$ 后，合理集合 $\mathcal{U}(r)$ 内仍存在使另一侧明显更快的流场 $\boldsymbol{u}^*$，继续加密网格也无法判断真实海洋接近 $\boldsymbol{u}_0$ 还是 $\boldsymbol{u}^*$。

第二，是否值得去观测，不取决于是否存在合理反例，而取决于该反例成真会让当前决策多损失多少小时（最坏后悔）以及观测能否触及关键切变区。由于旧粗图误差相关长度达数十千米，几乎任何任务都能找到某个让另一侧略快的合理流场；唯有直接求出所选一侧的最大可能航时劣势并与观测代价比较，才能避免过度反应。
